\section{Simulation studies}\label{sec:simulation}
\subsection{Design and competing analyses}
The experiments separate interval calibration for a known regression surface from sensitivity to changes in the dynamics of raw time series. This distinction permits coverage to be assessed against an exact coefficient in the first experiment and avoids assigning an unverified latent-spectrum target to finite-window responses in the second. Performance measures include global rejection probability, simultaneous coverage, average band width, and root mean squared error (RMSE), following the distinction between estimands and operating characteristics emphasized by \citet{Morris2019}. Each setting uses 1,000 Monte Carlo replications and 399 bootstrap draws per dataset. RMSE is the square root of the mean squared coefficient error, with averaging over both replications and grid points. Tests use a nominal level of 0.05 and bands use a nominal coverage of 0.95. Monte Carlo standard errors for rejection and coverage probabilities are at most about 0.016 and are about 0.007 near a probability of 0.05 or 0.95. Error bars use 95\% Wilson binomial intervals.

In the first experiment, the response is a smooth frequency--quantile surface with a subject random surface and serially correlated segment errors. There are two subject-level covariates and one segment-level covariate. The exposure coefficient is zero under the null and has two frequency-localized components under the alternative, with magnitude increasing away from the median quantile. Balanced designs use $N\in\{80,160\}$ subjects, $m\in\{4,12\}$ segments, and exposure coefficient magnitudes $\delta\in\{0,0.35\}$. An informative-count design uses $N=160$, $\delta=0$, and segment counts that depend on the exposure and a subject random effect. The complete generating equations are given in the appendix.

The proposed subject-weighted cluster analysis is compared with an analysis that treats all segments as independent and with a regression of subject-averaged responses on subject-averaged covariates. The latter uses the subject as its sampling unit but reduces the information available to estimate the segment-level covariate effect. The informative-count design also includes a cluster analysis that weights each segment equally. The same bootstrap construction is used for each method at its assumed independent sampling unit. A Gaussian multiplier procedure with fixed standard errors is examined separately in the appendix.

The second experiment generates Gaussian autoregressive (AR) and stochastic-volatility (SV) series, each with a subject-specific persistence parameter related to a binary exposure. The models have six segments per subject. Segment length is 512 observations in the main comparison, with a longer-window setting used to assess sensitivity. Quantile spectra are evaluated at $\tau\in\{0.1,0.3,0.5,0.7,0.9\}$ and 30 frequencies from 0.015 to 0.30 cycles per observation. The Gaussian smoothing bandwidth is 0.02 and $\epsilon=10^{-6}$. Both null settings and exposure effects of several magnitudes are included.

For Gaussian autoregressions, the exposure changes ordinary and quantile spectral dependence. For the volatility model, the marginal distribution and population ordinary spectrum are unchanged by the exposure, while the persistence of noncentral quantile events changes. We compare the full quantile regression, median-only spectral regression, ordinary log-power regression, and the independent-segment quantile analysis. The ordinary and median-only analyses are useful reference procedures for different features of the process, rather than interchangeable estimators of the full quantile surface.
